% Part: computability
% Chapter: computability-theory
% Section: non-comp-set

\documentclass[../../../include/open-logic-section]{subfiles}

\begin{document}

% SEGMENT OLP-0240-S01

\olfileid{cmp}{thy}{ncp}
\olsection{கணிக்கத்தகாத கணங்கள் உள்ளன}


ஒவ்வொரு கணிக்கத்தக்க கணமும் கணிக்கத்தக்கவாறு எண்ணிடத்தக்கது என்பதை
மேலே கண்டோம். இதன் மறுதலை உண்மையா? பொதுவாக அது உண்மையல்ல என்பதைப்
பின்வரும் தேற்றம் காட்டுகிறது.

\begin{thm}\ollabel{thm:K-0}
$K_0$ என்பது $\Setabs{\tuple{e, x}}{\cfind{e}(x) \fdefined}$ என்ற
கணம் ஆகட்டும். அப்போது $K_0$ கணிக்கத்தக்கவாறு எண்ணிடத்தக்கது;
ஆனால் கணிக்கத்தக்கது அல்ல.
\end{thm}

\begin{proof}
$K_0$ கணிக்கத்தக்கவாறு எண்ணிடத்தக்கது என்பதைக் காண, பின்வருமாறு
வரையறுக்கப்பட்ட சார்பு~$f$ இன் வரையறைப்புலமாக அது இருப்பதைக் கவனிக்க:
\[
f(z) = \umin{y}{(\len{z} = 2 \land T((z)_0, (z)_1, y))}.
\]
ஏனெனில், $\cfind{e}(x)$ வரையறுக்கப்பட்டிருந்தால்,
$f(\tuple{e, x})$ ஒரு நிற்கும் கணிப்புத் தொடரைக் கண்டுபிடிக்கிறது;
$\cfind{e}(x)$ வரையறுக்கப்படவில்லை எனில் $f(\tuple{e, x})$ உம்
வரையறுக்கப்படவில்லை; $z$ ஒரு சோடியைக்கூட குறிமுறையாக்கவில்லை எனில்,
$f(z)$ உம் வரையறுக்கப்படவில்லை.

$K_0$ கணிக்கத்தக்கது அல்ல என்பது நிற்றல் சிக்கலின்
தீர்மானிக்கவியலாமையே; காண்க \olref[hlt]{thm:halting-problem}.
\end{proof}

$K_0$ என்பது $\cfind{e}(x) \fdefined$ ஆகும் சோடிகள்
$\tuple{e,x}$ இன் கணம்; அதாவது $\cfind{e}$ ஆனது உள்ளீடு~$x$ இல்
வரையறுக்கப்பட்டு (நின்று) இருந்தால், எனில் மட்டுமே,
$\tuple{e,x} \in K_0$. எனவே அது ``நிற்கும் கணம்'' என்றும்
அழைக்கப்படுகிறது. $K = \Setabs{e}{\cfind{e}(e) \fdefined}$ என்பது
``தன்னிலேயே நிற்கும் கணம்.'' இது ஒரு நியமமான
தீர்மானிக்கவியலாத கணமாக அடிக்கடி பயன்படுத்தப்படுகிறது.

\begin{thm}\ollabel{thm:K}
தன்னிலேயே நிற்கும் கணம்
$K = \Setabs{e}{\cfind{e}(e) \fdefined}$ ஆனது !!{c.e.};
ஆனால் தீர்மானிக்கத்தக்கது அல்ல.
\end{thm}

\begin{proof}
  $K$ தீர்மானிக்கத்தக்கது, அதாவது அதன் சிறப்பியல்புச் சார்பு
  $\Char{K}$ கணிக்கத்தக்கது எனக் கொள்க.
  \[d(e) = \begin{cases}
  1 & \text{$\Char{K}(e) = 0$ எனில்}\\
  \fundefined & \text{இல்லையெனில்.}
  \end{cases}
  \]
  $k$ ஆனது~$d$ இன் சுட்டெண், அதாவது
  $d \simeq \cfind{k}$ ஆகட்டும். அப்போது
  $d(k) \simeq \cfind{k}(k)$. ஆனால் $d$ இன் வரையறையிலிருந்து,
  $\cfind{k}(k) \fundefined$ எனில், எனில் மட்டுமே,
  $d(k) \fdefined$ என்பது கிடைக்கிறது; இது முரண்பாடு.

  $K$ ஆனது $f(x) = \umin{y}{T(x,x,y)}$ இன் வரையறைப்புலம்;
  எனவே அது !!{c.e.}
\end{proof}

\end{document}
